\section{工具不是附属品}

LCQI 的工程证据不只在 C++ 源码里。smoke 脚本决定真实模型怎样下载、校验和运行；benchmark 程序决定性能口径；trace、ledger 和 serving probe 决定报告字段；CTest 决定哪些行为能长期回归。读者必须把这些文件当作工程合同的一部分。

\section{Benchmark、Trace、Ledger 与 Serving Probe}

\filepath{bench.cpp} 输出 int8 kernel benchmark。它服务性能观察，但必须和 max diff 对照一起看。

\lstinputlisting[language=C++,caption={src/bench.cpp}]{../../../cpu-volume-3/labs/linux_cpu_inference/src/bench.cpp}

\filepath{trace.cpp} 输出 reference decoder checkpoint。它用于定位错误边界，不应该被简化成只打印最终 token。

\lstinputlisting[language=C++,caption={src/trace.cpp}]{../../../cpu-volume-3/labs/linux_cpu_inference/src/trace.cpp}

\filepath{ledger.cpp} 输出 7B 规模账本。它是估算，不是实测吞吐；读者要区分 FLOPs、KV 容量、权重/KV 字节和 bandwidth-only tokens/s 上限。

\lstinputlisting[language=C++,caption={src/ledger.cpp}]{../../../cpu-volume-3/labs/linux_cpu_inference/src/ledger.cpp}

\filepath{serving\_probe.cpp} 固定短请求、RAG 请求、取消请求和超长请求。它把 TTFT、TPOT、拒绝和取消回收变成可观察字段。

\lstinputlisting[language=C++,caption={src/serving\_probe.cpp}]{../../../cpu-volume-3/labs/linux_cpu_inference/src/serving_probe.cpp}

\filepath{onednn\_bench.cpp} 是可选外部库对标入口。它只在环境安装 oneDNN 时构建，不能把不可用平台写成已验证性能。

\lstinputlisting[language=C++,caption={src/onednn\_bench.cpp}]{../../../cpu-volume-3/labs/linux_cpu_inference/src/onednn_bench.cpp}

\filepath{gguf\_q4\_bench.cpp} 是真实 GGUF 量化 tensor 的单算子 benchmark。它读取同一个 SmolLM2 Q4\_K\_M GGUF，选择真实 \code{Q4\_K} 矩阵，分别测 f32 解码基线、\code{Q4\_K x f32} 和 \code{Q4\_K x Q8\_K}。它输出的 \code{speedup} 只适用于这个单 tensor matvec，不能写成完整模型 tokens/s。

\lstinputlisting[language=C++,caption={src/gguf\_q4\_bench.cpp}]{../../../cpu-volume-3/labs/linux_cpu_inference/src/gguf_q4_bench.cpp}

\filepath{q4\_k\_avx2.cpp} 是本轮 Q4\_K 优化入口。它不是完整 llama.cpp 后端，只把 \code{Q4\_K x Q8\_K} 的 32 元素子块点积换成 AVX2 \code{maddubs} 归约。读这份源码时要同时看 fallback：环境变量 \code{LCQI\_Q4K\_BACKEND=scalar} 可以强制走标量路径，报告脚本用它生成同机同模型的 A/B 证据。

\lstinputlisting[language=C++,caption={src/q4\_k\_avx2.cpp}]{../../../cpu-volume-3/labs/linux_cpu_inference/src/q4_k_avx2.cpp}

\section{外部模型 Smoke 与对比脚本}

\filepath{run\_smollm2\_small\_smoke.py} 下载并校验 SmolLM2-135M-Instruct 的 GGUF 量化文件，构建固定 commit 的 \code{llama-simple-chat}，在 CPU 上运行一次真实开源 small 模型生成，并输出可复查报告。它是外部真实模型 smoke，不替代 LCQI tiny golden trace。

\lstinputlisting[language=Python,caption={tools/run\_smollm2\_small\_smoke.py}]{../../../cpu-volume-3/tools/run_smollm2_small_smoke.py}

\filepath{run\_smollm2\_q4\_k\_benchmark.py} 复用同一个 SmolLM2 GGUF 缓存，构建 LCQI 的 \code{lcqi\_gguf\_q4\_bench}，并可附带 llama.cpp \code{llama-bench} 端到端参考。它会先用 \code{LCQI\_Q4K\_BACKEND=scalar} 跑标量路径，再跑 auto 后端，因此报告里能直接看到 \code{avx2\_maddubs} 相对 scalar 的提升。报告要同时说明模型 SHA-256、tensor 名、shape、量化字节、f32 解码字节、误差、checksum 和成熟引擎参考，避免把不同层级数字混在一起。

\lstinputlisting[language=Python,caption={tools/run\_smollm2\_q4\_k\_benchmark.py}]{../../../cpu-volume-3/tools/run_smollm2_q4_k_benchmark.py}

\filepath{run\_gpt2\_smoke.py} 下载标准 GPT-2 的 \code{config.json}、\code{vocab.json}、\code{merges.txt} 和 \code{model.safetensors}，再调用 LCQI 自研 \code{lcqi\_gpt2} reference path 生成 token。它证明自研 safetensors、BPE、GPT-2 forward、KV cache decode 和 greedy generation 已连成闭环。

\lstinputlisting[language=Python,caption={tools/run\_gpt2\_smoke.py}]{../../../cpu-volume-3/tools/run_gpt2_smoke.py}

\filepath{run\_gpt2\_benchmark\_compare.py} 用同一个 GPT-2 F32 safetensors 模型分别运行 LCQI full-prefix 与 cached-KV 路径，并在本机存在 llama.cpp/SmolLM2 缓存时附带成熟引擎参考。它的报告说明算法差距、kernel 差距和模型口径差异。

\lstinputlisting[language=Python,caption={tools/run\_gpt2\_benchmark\_compare.py}]{../../../cpu-volume-3/tools/run_gpt2_benchmark_compare.py}

\filepath{run\_gpt2\_optimization\_ab.py} 专门比较 LCQI 自身改动前后。它从 baseline commit 建临时 worktree，baseline 和当前工作区都通过实践卷 CMake 入口构建 Release，再多轮运行同一个 GPT-2 F32 模型。报告中的 \code{median/min/max} 用来约束机器波动，\code{generated\_text} 用来确认优化前后输出没有漂移，\code{benchmark\_worker\_count} 用来确认 cached decoder 到底是串行、自动线程数还是固定线程数。

\lstinputlisting[language=Python,caption={tools/run\_gpt2\_optimization\_ab.py}]{../../../cpu-volume-3/tools/run_gpt2_optimization_ab.py}

\filepath{run\_gpt2\_hotspot\_profile.py} 专门收集 cached decode 内部热点。它构建 Release \code{lcqi\_gpt2}，用 \code{--profile-hotspots} 输出 \code{hotspot\_lm\_head\_pct}、\code{hotspot\_mlp\_fc\_pct}、\code{hotspot\_mlp\_projection\_pct}、\code{hotspot\_qkv\_projection\_pct}、\code{hotspot\_attention\_projection\_pct}、\code{hotspot\_attention\_pct} 和 \code{benchmark\_worker\_count}，再按 token 数和轮次计算中位数。这个脚本回答“优化应该打哪里”，不是回答“某次提交快了多少”。本轮线程优化就是靠它发现第一次 auto 阈值漏掉了 attention output projection，随后把行数和列数阈值一起纳入并行判定。

\lstinputlisting[language=Python,caption={tools/run\_gpt2\_hotspot\_profile.py}]{../../../cpu-volume-3/tools/run_gpt2_hotspot_profile.py}

\filepath{run\_lcqi\_benchmark.sh} 是本地 benchmark 收集脚本。它适合在固定机器上反复运行，把同一组 shape 和 backend 的结果落到报告文件。

\lstinputlisting[language=bash,caption={tools/run\_lcqi\_benchmark.sh}]{../../../cpu-volume-3/tools/run_lcqi_benchmark.sh}

\section{完整测试}

\filepath{lcqi\_tests.cpp} 把 safetensors、tokenizer、tiny MLP、int8 kernel、reference decoder、GPT-2 tiny path、trace、ledger 和 serving probe 串成验收网。GPT-2 优化不是只看速度：测试会把 full-prefix logits、旧 cached logits、优化 decoder logits、强制两 worker 的并行 decoder logits 和 logits-free greedy token 放在同一个 tiny prompt 下对齐，确认 predicted token、logits、KV filled token 数和 generated ids 不漂移。新增优化或 shape 时，必须先扩展这里的 golden 和边界样例。

\lstinputlisting[language=C++,caption={tests/lcqi\_tests.cpp}]{../../../cpu-volume-3/labs/linux_cpu_inference/tests/lcqi_tests.cpp}
